% GENERATED FILE. Edit canonical lessons, then run npm run generate:books.
% canonical-source-hash: fnv1a64:03d6a74b00e2f24e
% canonical-lessons: GE-C338-vermutung, GE-C338-annahme, GE-C338-anscheinend, GE-C338-offensichtlich, GE-C338-ahnen

\chapter{Vermutungen}
\label{ch:ge-vermutungen}
\hlchaptermodality{338}

\begin{chapteropening}
\textbf{By the end of this chapter:} \emph{I can make a guess, state an assumption, and say how things seem or plainly are.}

Complete the last lesson of chapter 338: I can make a guess, state an assumption, and say how things seem or plainly are.
\end{chapteropening}

\section[die Vermutung]{die Vermutung — the guess, the supposition}

\label{lesson:GE-C338-vermutung}



\begin{quote}
Before the new one: say the German for to have doubts, then the German for to doubt.
\end{quote}

\subsection*{You'll want to know: die Vermutung}
\textbf{die Vermutung} — \textquotedblleft{}the guess, the supposition\textquotedblright{}.

\begin{quote}
Meine Vermutung war richtig: Er hatte den Schlüssel im Auto vergessen.
\end{quote}

\begin{quote}
\emph{My guess was right: he had left the key in the car.}
\end{quote}

Plural: \emph{die Vermutungen}. The noun of \emph{vermuten}.

\begin{cousinweb}
\emph{Vermutung} is the noun of \emph{vermuten} \textquotedblleft{}to suppose\textquotedblright{}, which is usually traced to \emph{muten}, from \emph{Mut} in its old sense \textquotedblleft{}mind, feeling\textquotedblright{}. EtymDB has no entry for the verb, so the account is hedged.
\end{cousinweb}

\begin{grammarlens}[title={What you've built}]
One more word for telling, explaining and describing at B1.
\end{grammarlens}

\subsection*{Your turn}
\begin{itemize}
\raggedright
  \item \emph{Say it:} \emph{die Vermutung}
  \item \emph{Say it:} \emph{die Vermutung}, once more
  \item \emph{Say it:} a sentence of your own with \emph{die Vermutung}
  \item \emph{Recall:} say the German for to have doubts, then the German for to doubt, then say \emph{die Vermutung} again
\end{itemize}

\begin{tcolorbox}[breakable,colback=teal!4,colframe=teal!35!black,title={Before you move on}]
What does die Vermutung mean? (\textquotedblleft{}The guess, the supposition\textquotedblright{}.) Say it once more.
\end{tcolorbox}
\section[die Annahme]{die Annahme — the assumption; the acceptance}

\label{lesson:GE-C338-annahme}



\begin{quote}
Before the new one: say the German for to doubt, then the German for the guess.
\end{quote}

\subsection*{You'll want to know: die Annahme}
\textbf{die Annahme} — \textquotedblleft{}the assumption; the acceptance\textquotedblright{}.

\begin{quote}
Wir sind von der falschen Annahme ausgegangen, dass das Museum montags offen ist.
\end{quote}

\begin{quote}
\emph{We started from the wrong assumption that the museum is open on Mondays.}
\end{quote}

Plural: \emph{die Annahmen}. The noun of \emph{annehmen}: \emph{die Annahme eines Pakets} — accepting a parcel.

\begin{cousinweb}
\emph{Annahme} is the noun of \emph{annehmen}, built on \emph{nehmen} \textquotedblleft{}to take\textquotedblright{}, inherited from Old High German \emph{neman} and Proto-Germanic \emph{nemaną} (EtymDB): something you take as given.
\end{cousinweb}

\begin{grammarlens}[title={What you've built}]
One more word for telling, explaining and describing at B1.
\end{grammarlens}

\subsection*{Your turn}
\begin{itemize}
\raggedright
  \item \emph{Say it:} \emph{die Annahme}
  \item \emph{Say it:} \emph{die Annahme}, once more
  \item \emph{Say it:} a sentence of your own with \emph{die Annahme}
  \item \emph{Recall:} say the German for to doubt, then the German for the guess, then say \emph{die Annahme} again
\end{itemize}

\begin{tcolorbox}[breakable,colback=teal!4,colframe=teal!35!black,title={Before you move on}]
What does die Annahme mean? (\textquotedblleft{}The assumption; the acceptance\textquotedblright{}.) Say it once more.
\end{tcolorbox}
\section[anscheinend]{anscheinend — apparently, it seems}

\label{lesson:GE-C338-anscheinend}



\begin{quote}
Before the new one: say the German for the guess, then the German for the assumption.
\end{quote}

\subsection*{You'll want to know: anscheinend}
\textbf{anscheinend} — \textquotedblleft{}apparently, it seems\textquotedblright{}.

\begin{quote}
Anscheinend ist niemand zu Hause: Alle Fenster sind dunkel.
\end{quote}

\begin{quote}
\emph{Apparently nobody is at home: all the windows are dark.}
\end{quote}

It says how things look, and they probably are so. \emph{Scheinbar} says they only look so.

\begin{cousinweb}
\emph{anscheinend} is the present participle of \emph{anscheinen} \textquotedblleft{}to appear to\textquotedblright{}, built on \emph{scheinen} \textquotedblleft{}to shine, to seem\textquotedblright{}, inherited from Old High German \emph{skīnan} and Proto-Germanic \emph{skīnaną} (EtymDB), a cousin of English \emph{shine}.
\end{cousinweb}

\begin{grammarlens}[title={What you've built}]
One more word for telling, explaining and describing at B1.
\end{grammarlens}

\subsection*{Your turn}
\begin{itemize}
\raggedright
  \item \emph{Say it:} \emph{anscheinend}
  \item \emph{Say it:} \emph{anscheinend}, once more
  \item \emph{Say it:} a sentence of your own with \emph{anscheinend}
  \item \emph{Recall:} say the German for the guess, then the German for the assumption, then say \emph{anscheinend} again
\end{itemize}

\begin{tcolorbox}[breakable,colback=teal!4,colframe=teal!35!black,title={Before you move on}]
What does anscheinend mean? (\textquotedblleft{}Apparently, it seems\textquotedblright{}.) Say it once more.
\end{tcolorbox}
\section[offensichtlich]{offensichtlich — obviously, evident}

\label{lesson:GE-C338-offensichtlich}



\begin{quote}
Before the new one: say the German for the assumption, then the German for apparently.
\end{quote}

\subsection*{You'll want to know: offensichtlich}
\textbf{offensichtlich} — \textquotedblleft{}obviously, evident\textquotedblright{}.

\begin{quote}
Offensichtlich hat er den Brief nicht gelesen.
\end{quote}

\begin{quote}
\emph{Obviously he has not read the letter.}
\end{quote}

An adjective and an adverb: \emph{ein offensichtlicher Fehler} — an obvious mistake.

\begin{cousinweb}
\emph{offensichtlich} joins \emph{offen} \textquotedblleft{}open\textquotedblright{} to \emph{Sicht} \textquotedblleft{}sight\textquotedblright{}, from the family of \emph{sehen} \textquotedblleft{}to see\textquotedblright{}, Old High German \emph{sehan} (EtymDB), with \emph{-lich}: open to sight.
\end{cousinweb}

\begin{grammarlens}[title={What you've built}]
One more word for telling, explaining and describing at B1.
\end{grammarlens}

\subsection*{Your turn}
\begin{itemize}
\raggedright
  \item \emph{Say it:} \emph{offensichtlich}
  \item \emph{Say it:} \emph{offensichtlich}, once more
  \item \emph{Say it:} a sentence of your own with \emph{offensichtlich}
  \item \emph{Recall:} say the German for the assumption, then the German for apparently, then say \emph{offensichtlich} again
\end{itemize}

\begin{tcolorbox}[breakable,colback=teal!4,colframe=teal!35!black,title={Before you move on}]
What does offensichtlich mean? (\textquotedblleft{}Obviously, evident\textquotedblright{}.) Say it once more.
\end{tcolorbox}
\section[ahnen]{ahnen — to suspect, to have a feeling}

\label{lesson:GE-C338-ahnen}



\begin{quote}
Before the new one: say the German for apparently, then the German for obviously.
\end{quote}

\subsection*{You'll want to know: ahnen}
\textbf{ahnen} — \textquotedblleft{}to suspect, to have a feeling\textquotedblright{}.

\begin{quote}
Ich habe es schon geahnt: Der Zug hat Verspätung.
\end{quote}

\begin{quote}
\emph{I already had a feeling: the train is late.}
\end{quote}

\emph{Keine Ahnung} — no idea — belongs to the same family.

\begin{cousinweb}
\emph{ahnen} is inherited from Middle High German \emph{anen} (EtymDB). It is usually explained as a verb made from \emph{an} \textquotedblleft{}on\textquotedblright{}: a feeling comes upon you.
\end{cousinweb}

\begin{grammarlens}[title={What you've built}]
One more word for telling, explaining and describing at B1.
\end{grammarlens}

\subsection*{Your turn}
\begin{itemize}
\raggedright
  \item \emph{Say it:} \emph{ahnen}
  \item \emph{Say it:} \emph{ahnen}, once more
  \item \emph{Say it:} a sentence of your own with \emph{ahnen}
  \item \emph{Recall:} say the German for apparently, then the German for obviously, then say \emph{ahnen} again
\end{itemize}

\begin{tcolorbox}[breakable,colback=teal!4,colframe=teal!35!black,title={Before you move on}]
What does ahnen mean? (\textquotedblleft{}To suspect, to have a feeling\textquotedblright{}.) Say it once more.
\end{tcolorbox}
